%%%%%%%%%%%%%%%%%%%%%%%%%%%%%%%%%%%%%%%%%%%%%%%%%%%%%%%%%
%%%%%%%%%%%%%%%%%%         APENDICES             %%%%%%%%%%%%%%%%%%%%%%%
%%%%%%%%%%%%%%%%%%%%%%%%%%%%%%%%%%%%%%%%%%%%%%%%%%%%%%%%%


\chapter{Producción científica}
\label{chap:apendices}

\chapterintro{En esta sección se recopila la producción científica generada durante el período de tesis. Se marca la diferencia entre los tres artículos que componen esta tesis y entre los artículos y/o participaciones en congresos generados a raíz del desarrollo de la tesis}

\section{Compendio de tres artículos JCR (WOS)}
\begin{itemize}

    \item Loaiza, L. I. S., \textbf{Fernández-Edreira, D.}, Liñares-Blanco, J., Cepeda, A., Cardelle-Cobas, A., and Fernandez-Lozano, C. (2025). Fecal microbiome analysis in patients with metabolic syndrome and type 2 diabetes. PeerJ, 13, e19108. Q2, 2.4 IF. 1 cita (Google Scholar). \url{https://doi.org/10.7717/peerj.19108}
    
    \item \textbf{Fernández-Edreira, D.}, Liñares-Blanco, J., and Fernandez-Lozano, C. (2021). Machine Learning analysis of the human infant gut microbiome identifies influential species in type 1 diabetes. Expert Systems with Applications, 185, 115648. Q1, 8.665 IF. 48 citas (Google Scholar). \url{https://doi.org/10.1016/j.eswa.2021.115648}

    \item \textbf{Fernández-Edreira, D.}, Liñares-Blanco, J., Patricia, V., and Fernandez-Lozano, C. (2024). VIBES: A consensus subtyping of the vaginal microbiota reveals novel classification criteria. Computational and Structural Biotechnology Journal, 23, 148-156. Q2, 4.1 IF. 4 citas (Google Scholar). \url{https://doi.org/10.1016/j.csbj.2023.11.050}

    
\end{itemize}

\section{Otros artículos JCR (WOS) (1)}

\begin{itemize}

    \item Golob, J. L., Oskotsky, T. T., Tang, A. S., Roldan, A., et al. (2024). Microbiome preterm birth DREAM challenge: Crowdsourcing machine learning approaches to advance preterm birth research. Cell Reports Medicine, 5(1). Q1, 10.6 IF. \url{https://doi.org/10.1016/j.xcrm.2023.101350}

\end{itemize}

\section{Otros artículos (1)}

\begin{itemize}

    \item \textbf{Fernández-Edreira, D.}, Liñares-Blanco, J., Galdo, B. C., García, V. M., et al. (2021). La Inteligencia Artificial como herramienta para la gestión y explotación de datos, informaciones y conocimientos biomédicos en entornos “Big Data” en la nube. I+ S: Revista de la Sociedad Española de Informática y Salud, 143, 30-39. \url{https://oa.upm.es/81987/1/MAOJO_ART_JOUR_113.pdf}

\end{itemize}

\section{Participación en congresos nacionales e internacionales (3)}

\begin{itemize}

    \item Filipe, H. A., Costa, A. M., Moreira, A., Miguel, S. M., et al. (2025). RePo-SUDOE: A Transnational Network for Drug Repurposing in the SUDOE Space. Drug Repurposing, 2(1), 20250004. \url{https://doi.org/10.58647/DRUGREPO.25.1.0004}

    \item Quiza, R. M., \textbf{Fernández-Edreira, D.}, and Fernandez-Lozano, C. (2025). Facilitando el aprendizaje en bioinformática y reposicionamiento de fármacos mediante una plataforma web intuitiva enfocada en el análisis de resultados. Actas de las JENUI, 10, 311-316. \url{https://aenui.org/actas/pdf/JENUI_2025_035.pdf}

    \item \textbf{Fernández-Edreira, D.}, Liñares-Blanco, J., and Fernandez-Lozano, C. (2024). Cross Sequencing Integration of Compositional Microbiome Data in Cancer. In International Meeting on Computational Intelligence Methods for Bioinformatics and Biostatistics (pp. 69-78). Cham: Springer Nature Switzerland. \url{https://doi.org/10.1007/978-3-031-89704-7_6}

\end{itemize}




